% ============================================================
% Section 3: Numerical Schemes
% ============================================================

\section{Numerical Schemes and Implementation}
\pagecheck{2--3}

\instructionplaceholder{Implementation logic}{Document how the discrete FEM equations are implemented in Python or MATLAB. The reader should be able to understand solver workflow, data flow, and time integration choices from this section alone. You may use the lecture notebooks as a starting point; more efficient implementations will be valued positively.}

\subsection{Time integration and solver settings}
\instructionplaceholder{Required evidence}{State time-stepping methods, stability constraints (e.g.\ Courant number, critical time step), solver tolerances, and convergence criteria. Explain why the selected scheme is appropriate for the structural dynamics of the FOWT system, taking into account the relevant natural frequencies identified in Section 4.}

\subsection{Algorithm design and code structure}
\instructionplaceholder{Required evidence}{Describe the computational pipeline, key functions/modules, and reproducibility setup. Include the sequence of assembly, load application, time stepping, and post-processing. Highlight efficiency choices and any extensions beyond the lecture notebooks.}

\subsection{Implementation ownership checklist}
\instructionplaceholder{Required evidence}{For each major solver component, report: selected approach, one alternative considered, and why the final choice was implemented. Detailed convergence, stability, and modal verification evidence belongs in Section 4.}
